% Generated from locale/tr/content/first-order-logic/syntax-and-semantics/introduction.tex; do not hand-edit.


\olfileid[tr]{fol}{syn}{int}

\olsection{Giriş}

Birinci dereceden mantığın kuramını ve metakuramını geliştirmek için önce
ifadelerinin sözdizimini ve semantiğini tanımlamamız gerekir. Birinci dereceden
mantığın ifadeleri terimler ve !!{formula}s{}dir. Terimler !!{variable}s,
!!{constant}s ve !!{function}s kullanılarak oluşturulur. !!^{formula}s ise
!!{predicate}s ile terimlerden (bunlar en küçük, ``atomik'' !!{formula}s{}i
oluşturur), ardından da atomik !!{formula}s{}den mantıksal bağlaçlar ve
niceleyiciler kullanılarak oluşturulur. Kuruluş kurallarını belirlemenin pek
çok yolu vardır; burada yalnızca bunlardan birini veriyoruz. Başka sistemler
başka simgeler seçer, başka bağlaç kümelerini ilkel kabul eder ve parantezleri
başka biçimde kullanır (hatta Leh gösterimi denen gösterimde olduğu gibi hiç
kullanmaz). Bununla birlikte bütün yaklaşımların ortak yanı, kuruluş
kurallarının terimler ve !!{formula}s kümesini \emph{tümevarımsal olarak}
tanımlamasıdır. Kurallar doğru kurulmuşsa her ifade, kuruluş kurallarına göre
özünde yalnızca tek bir biçimde elde edilebilir. \emph{Biricik okunabilir}
ifadeler üreten tümevarımsal tanım, aynı yöntemle---tümevarımsal tanımla---bu
ifadelere anlam verebilmemizi sağlar.

İfadelere anlam vermek semantiğin alanıdır. Semantiğin merkezî kavramı,
!!a{structure} içinde sağlanmadır. !!^a{structure}, dilin yapı taşlarına anlam
verir: !!a{domain}, boş olmayan bir nesneler kümesidir. Niceleyiciler bu
!!{domain} üzerinde değişecek biçimde yorumlanır; !!{constant}s bu kümenin
elemanlarına, !!{function}s !!{domain}{}den yine kendisine giden fonksiyonlara ve
!!{predicate}s !!{domain} üzerindeki bağıntılara atanır. !!{domain} ile temel
söz dağarcığına yapılan atamalar birlikte !!a{structure} oluşturur.
!!^{variable}s, !!{formula}s içinde bulunabilir; bunlara da semantik vermek
için !!{domain}{}in !!{element}s{}ını atamamız gerekir---bu bir değişken
atamasıdır. Son olarak sağlanma bağıntısı bütün bunları bir araya getirir.
!!^a{formula}, bir $s$ değişken atamasına göre !!a{structure}~$\Struct{M}$
içinde sağlanabilir; bunu $\Sat{M}{!A}[s]$ biçiminde yazarız. Bu bağıntı da
$!A$'nın yapısı üzerinde tümevarımla tanımlanır: Örneğin $!A \land !B$'nin
sağlanması, mantıksal bağlaçların doğruluk tabloları kullanılarak $!A$ ile
~$!B$'nin sağlanması (ya da sağlanmaması) cinsinden tanımlanır. Sonra,
!!{formula}~$!A$ serbest değişken içermeyen !!a{sentence} olduğunda değişken
atamasının önemsiz olduğu ortaya çıkar; böylece !!{sentence}s{}in
!!{structure}s içinde doğrudan sağlanmasından (ya da sağlanmamasından) söz
edebiliriz.

Tümceler için $\Sat{M}{!A}$ sağlanma bağıntısına dayanarak geçerlilik,
mantıksal gerektirme ve sağlanabilirliğin temel semantik kavramlarını
tanımlayabiliriz. Bir tümce, her yapı onu sağlıyorsa geçerlidir,
$\Entails !A$. Bir !!{sentence}s kümesi tarafından mantıksal olarak
gerektirilmesi, $\Gamma \Entails !A$, $\Gamma$'daki bütün !!{sentence}s{}i
sağlayan her !!{structure}{}nın $!A$'yı da sağlaması demektir. Bir tümceler
kümesi de, bir !!{structure} içindeki bütün !!{sentence}s{}i aynı anda
sağlıyorsa sağlanabilirdir. !!{formula}s tümevarımsal olarak, sağlanma da
!!{formula}s{}in yapısı üzerinde tümevarımla tanımlandığı için semantiğimizin
özelliklerini ispatlamak ve tanımlanan semantik kavramları ilişkilendirmek
üzere tümevarımı kullanabiliriz.




